\documentclass[11pt]{article}
\usepackage[letterpaper,margin=0.88in]{geometry}
\usepackage[T1]{fontenc}
\usepackage{lmodern,amsmath,amssymb,amsthm,mathtools}
\usepackage{booktabs,longtable,array,microtype,xcolor,enumitem,fancyhdr}
\usepackage[colorlinks=true,linkcolor=blue!45!black,urlcolor=blue!50!black,citecolor=blue!45!black]{hyperref}
\usepackage{xurl}
\hypersetup{pdftitle={A Direct Integer Equation for the Consistency of ZFC},pdfauthor={ChatGPT, OpenAI; prepared for Vladimir Reshetnikov}}
\newtheorem{theorem}{Theorem}[section]
\newtheorem{lemma}[theorem]{Lemma}
\newtheorem{proposition}[theorem]{Proposition}
\newtheorem{corollary}[theorem]{Corollary}
\theoremstyle{definition}\newtheorem{definition}[theorem]{Definition}
\theoremstyle{remark}\newtheorem{remark}[theorem]{Remark}
\newcommand{\Z}{\mathbb Z}\newcommand{\N}{\mathbb N}
\newcommand{\Con}{\operatorname{Con}}\newcommand{\ZFC}{\mathrm{ZFC}}
\newcommand{\Refut}{\operatorname{Ref}}\newcommand{\degp}{\operatorname{deg}}
\newcommand{\abs}[1]{\left|#1\right|}\newcommand{\sqnorm}[1]{\left\lVert#1\right\rVert}
\setlength{\parskip}{4pt}\setlength{\parindent}{0pt}
\setlist{nosep,leftmargin=1.6em}
\pagestyle{fancy}\fancyhf{}\fancyhead[L]{\small Direct Diophantine consistency}
\fancyhead[R]{\small ProveIt research contribution}\fancyfoot[C]{\thepage}
\setlength{\headheight}{14pt}
\begin{document}
\hypersetup{pageanchor=false}
\begin{titlepage}
\vspace*{0.35in}
{\large COMPUTABILITY / DIOPHANTINE CONSISTENCY\par}
\vspace{0.32in}
{\Huge\bfseries A direct integer equation\par for the consistency of ZFC\par}
\vspace{0.3in}
{\Large Joint Pell certificates and exact integral elimination\par}
\vspace{0.45in}
{\large Prepared for Vladimir Reshetnikov and the ProveIt project\par}
{\large 25 October 2026\par}
\vspace{0.25in}
\begin{center}
\fbox{\begin{minipage}{0.86\textwidth}
\centering\vspace{9pt}
\textbf{Principal delivered equation}\par\vspace{5pt}
$P_{85}(y_0,\ldots,y_{84})=\displaystyle\sum_{j=0}^{23}r_j(\boldsymbol y)^2=0$\par\vspace{5pt}
85 unrestricted integer unknowns; exact degree 16,964;\par
1,788 live binary arithmetic operations.\par\vspace{7pt}
$\displaystyle\ZFC\vdash
 \bigl(\neg\exists\boldsymbol y\in\Z^{85}\ P_{85}(\boldsymbol y)=0\bigr)
 \longleftrightarrow\Con(\ZFC).$
\vspace{9pt}
\end{minipage}}
\end{center}
\vspace{0.3in}
This statement uses the explicitly identified inherited C222/ZF/ZFC
calibration. The new arithmetic composition and coordinate eliminations
have ordinary mathematical proofs and exact finite certificates.
The appendices specify every coefficient and substitution of the principal
polynomial. No choice of coefficient depends on the truth of consistency.
\vfill
\textbf{Research status.} AI-assisted research with separately assigned
internal mathematical and artifact reviews. This is a concrete improvement
within the inspected direct-calculus family, not a global simplicity
minimum or a claim of worldwide priority. The archive also contains
alternative variable/degree tradeoffs. It supplies no new end-to-end
Lean or Rocq compilation and proves neither consistency nor inconsistency.
\end{titlepage}
\hypersetup{pageanchor=true}

\begin{abstract}
We construct a fixed polynomial over the integers whose integer
unsolvability is equivalent, provably in ZFC, to syntactic consistency of
ZFC. The principal polynomial has 85 unrestricted integer unknowns, exact
total degree 16,964, and a division-free evaluator with 1,788 live binary
operations. Its 24 residual squares and all polynomial reconstruction maps
are specified explicitly. The construction combines the cofinal CRT
anchor of ProveIt package 173 with the joint signed-index and power
certificate of package 176. Eleven old certificate coordinates are
replaced by one; 124 further coordinates are eliminated using acyclic
equations with unit coefficients. Each of these 124 coordinate
reconstructions is an integer polynomial.
The new composition retains the numerical C222 proof calculus and has no
K215 simulation, Jones polynomial, or large encoded-machine coefficient
tuple. We separate the unbounded number-theoretic and logical proof from
finite algebra verification, give exact degree and variable-dependence
certificates, and record the larger literal coefficients introduced by
the CRT-anchor parent. The result improves selected inspected direct
variable/degree/circuit bounds; equations in fewer variables already exist
in other regimes. A restricted investigation of additional Pell centers
is included to explain why a plausible competing route did not improve
its parent degree.
\end{abstract}

\setcounter{tocdepth}{1}
\tableofcontents
\newpage
\section{Result, scope, and comparison}
\label{sec:result}

The task is to give one actual polynomial, with all of its integer
coefficients determined, for which ZFC proves an equivalence between
integer unsolvability and its own syntactic consistency. Merely invoking
the existence part of the DPRM theorem would not specify such a
polynomial. A bounded family of consistency tests would also not meet
the task. This report gives a finite, independently replayable polynomial
definition and identifies the precise inherited theorem that supplies
its logical meaning.

There is no single established ordering of ``simplest.'' Number of
unknowns, degree, coefficient height, description length, evaluation
cost, and proof complexity can conflict. Our principal choice gives a
substantial improvement of an inspected direct proof-calculus family.
It is not the fewest-variable representation in the entire repository:
the repository already contains much smaller variable counts in
different degree and coefficient regimes. Literature on universal pairs
likewise treats variable count and degree jointly \cite{bayer2026}.

\begin{theorem}[Principal consistency equation]
\label{thm:main}
Let $P_{85}\in\Z[y_0,\ldots,y_{84}]$ be the polynomial specified in
Sections~\ref{sec:literal} and \ref{sec:elimination} and
Appendices~\ref{app:free}--\ref{app:source}.
Using the inherited calibration in Section~\ref{sec:logic}, ZFC proves
\[
 \exists\boldsymbol y\in\Z^{85}\ P_{85}(\boldsymbol y)=0
 \quad\longleftrightarrow\quad
 \exists p\in\N\ \Refut_{\ZFC}(p).
\]
Consequently ZFC proves
\[
 \neg\exists\boldsymbol y\in\Z^{85}\ P_{85}(\boldsymbol y)=0
 \quad\longleftrightarrow\quad\Con(\ZFC).
\]
The polynomial has exact total degree $16\,964$, depends on each of its
85 variables, and is a sum of 24 polynomial squares. Its supplied
division-free evaluator has 1,788 live binary additions, subtractions,
and multiplications.
\end{theorem}

The proof concerns the standard syntactic statement: no finite proof of
falsity in a fixed effective presentation of ZFC. It does not assume
consistency, infer the truth of arbitrary ZFC theorems, or require a
transitive model. All occurrences of ``integer'' in the delivered
equations mean an unrestricted integer. Positivity assertions below are
conclusions of retained equations or hypotheses of stated lemmas.

\subsection{Measured alternatives}
\begin{table}[htbp]
\centering
\caption{Delivered integer-polynomial tradeoffs. All coordinates range over $\Z$.}
\label{tab:metrics}
\small
\begin{tabular}{rrrrr}
\toprule
Variables & Exact degree & Residuals & Operations & $H$\\
\midrule
84 & 2\,097\,732 & 23 & 1\,775 & 7\,342\,557\\
\textbf{85} & \textbf{16\,964} & \textbf{24} & \textbf{1\,788} & \textbf{59\,805}\\
86 & 1\,092 & 25 & 1\,797 & 5\,341\\
87 & 956 & 26 & 1\,804 & 5\,093\\
88 & 444 & 27 & 1\,807 & 2\,149\\
90 & 256 & 29 & 1\,819 & 1\,297\\
100 & 64 & 39 & 1\,872 & 390\\
109 & 32 & 48 & 1\,895 & 308\\
\bottomrule
\end{tabular}
\end{table}

All degrees in this table are exact, not just estimates propagated
through a circuit. ``Operations'' counts live binary $+,-,\times$
instructions; a square and multiplication by a fixed integer each cost
one operation. This is an algebraic operation count, not a bit-cost
bound. The coefficient exponent $H$ certifies
$\max_{\alpha}|[\boldsymbol y^{\alpha}]P|\leq 2^H$; it is an upper
bound on expanded coefficients, not an exact height.

The comparison snapshot is ProveIt commit
\texttt{ecb9a4b4dda617512952525d50f17343700aac52}. Package~177
reports a direct 87-integer polynomial of degree $8\,388\,612$ with
2,447 operations, a 94-integer polynomial of degree 256 with 2,468
operations, a 104-integer polynomial of degree 64 with 2,451 operations,
and a 111-integer polynomial of degree 32 with 2,502 operations
\cite{p177}. Our 85-, 90-, 100-, and 109-integer profiles improve the
corresponding inspected dimension/degree/circuit comparisons.

That statement has an explicit cost qualification. The largest absolute
integer literal in the new supplied circuits is
$1\,099\,511\,627\,780=2^{40}+4$; package~177's corresponding direct
circuits use literals of magnitude at most 256. Large literal
coefficients arise in the inherited frame because the first requested
power has index $2^{20}t$. Thus the new equations do not dominate every
coefficient-sensitive simplicity measure. Neither actual expanded term
counts nor exact expanded coefficient heights are claimed for the
high-degree equations.

Two ZIPs were present in \texttt{docs/incoming} at the snapshot:
\path{ConZFC_Mixed_Base_Pooling_2026-10-25.zip} and
\path{ConZFC_Affine_Horizons_Centered_Certificates_2026-10-26.zip}.
Their binary contents were not accessible through the available
repository reader. They are outside the novelty comparison. The
mathematical theorem and literal equations do not depend on a novelty
claim about these unread submissions.

\subsection{What is new}
The generic ingredients are credited rather than renamed as new
theorems. Package~173 supplies the cofinal CRT index used as a Pell
anchor and the affine growth horizon. Package~176 proves joint
authentication of signed indices and powers. Acyclic elimination with
unit coefficients already appears in the earlier direct constructions,
including package~152. The contribution here is their explicit compatible
composition, the source-specific interface proof, and the resulting
frozen integer polynomials and certificates. In particular, the
147 retained quadratic rows, all sixteen power requests, and the actual
signed frame are checked together; the composition is not justified by
matching variable counts alone.

\section{The calibrated arithmetic source}
\label{sec:source}

The immediate source is package~173's
\texttt{source/anchor219/equation.json}, preserved as identical literal
polynomial data in \texttt{source/anchor219.json}. The working copy has
one additional trailing newline; the exact upstream bytes and their
authenticated Git blob identity are retained separately. It has 219 integer coordinates
$v_0,\ldots,v_{218}$ and 158 quadratic residuals $R_0,\ldots,R_{157}$.
The source profile enforces, in its arithmetic interpretation,
\[
 t=16+4d^2,\qquad K=2^{20}t,\qquad F=t^K,\qquad C=2v_{16}.
\]
These variable powers describe the interpretation of the polynomial
certificate. They are not operations in the final polynomial. In the
literal source $t=v_0$, $F=v_2$, and the shared Pell parameter is
$A=v_{137}+1$.

The source binds a numerical proof calculus directly to arithmetic.
Finite derivations are represented by local constructor and rule
conditions; a generalized binomial/CRT encoding enforces these local
conditions, and a Horn-flow argument recovers a finite refutation.
There is no intermediate Turing-machine computation in this branch.
The calibration is inherited, and its boundaries are stated in
Section~\ref{sec:logic}.

\subsection{Why the affine growth index is sufficient}
The deleted derived-label range of package~173 is replaced by a uniform
constructor bound. Its finite constructor-depth ledger uses
$\lambda=2^{20}$ and gives local positive labels below $t^{\lambda-1}$.
Write $L=t^{\lambda}$. The growth request $F=t^{\lambda t}$ then gives
$F=L^t$. The boundary case matters: an argument requiring $F>L^t$
would not suffice for this source.

The relevant large-evaluation lemma says that, for positive labels at
most $L$, $L\geq t$, and $F\geq L^t$, equality of two products of factors
$1+Fj$, with at most $t$ factors on one side and at most $2t+1$ on the
other, forces equality of their multisets of labels. First compare the
number of factors: an extra factor produces a value exceeding the
other side because $F+1>L^t$. Once both sides have $m\leq t$ factors,
their elementary symmetric coefficients lie in $[1,F]$ in positive
degrees. Differences lie strictly between $-F$ and $F$. Reduction at
the least nonzero coefficient modulo $F$ forbids a nonzero difference
polynomial from vanishing at $F$. Factorization over $\mathbb Q$ then
recovers the labels. Package~173 supplies the full boundary-valid
proof and the constructor-depth certificate \cite{p173}.

The other useful source change is cofinality. The actual CRT index $C$
can be chosen positive, even, and above all other requested indices
inside its prescribed residue class. The CRT modulus is odd, so the
allowed even representatives form an unbounded progression. The radix
and Pell frame are chosen afterward. Consequently $C$ can serve as the
semantic anchor without an unjustified bound on a previously fixed
representative. This ordering is retained unchanged here.

\section{A literal joint polynomial certificate}
\label{sec:literal}

Delete the following source coordinates and rows:
\begin{align*}
 D_v&=\{144,145,146,147,148,149,204,214,215,216,217\},\\
 D_R&=\{102,103,104,105,106,107,147,153,154,155,156\}.
\end{align*}
Every coordinate in $D_v$ occurs only in rows in $D_R$. Retain the other
208 coordinates and append $v_{219}=z$. The resulting 209-coordinate
system keeps the 147 rows $R_j$ for $j\notin D_R$ and replaces all eleven
deleted rows by the single residual defined below.

Let the five bases be
\[
 n_0=v_0,\quad n_1=v_3,\quad n_2=v_{189},\quad
 n_3=v_{100}-1,\quad n_4=v_{100}.
\]
For each row of Table~\ref{tab:requests}, $m_i$ is the requested output,
$k_i$ the requested exponent, and $(D_i,Y_i,Z_i)$ the stored Pell data.
The retained gates give $Z_i=AY_i$, and the actual abscissa is
$X_i=D_i+Z_i$. Affine shifts in the table are part of the exact
definition; they must not be dropped.

\begin{table}[htbp]
\centering
\caption{The sixteen literal power requests. The base in group $j$ is $n_j$.}
\label{tab:requests}
\begin{tabular}{rrrrrrr}
\toprule
$i$ & Group $j$ & $m_i$ & $k_i$ & $D_i$ & $Y_i$ & $Z_i$\\
\midrule
0 & 0 & $v_2$ & $2^{20}v_0$ & $v_{186}$ & $v_{187}$ & $v_{188}$\\
1 & 1 & $v_{17}$ & $v_1$ & $v_{177}$ & $v_{178}$ & $v_{179}+1$\\
2 & 2 & $v_{100}-1$ & $v_1$ & $v_{177}$ & $v_{178}$ & $v_{179}+1$\\
3 & 4 & $v_{101}$ & $v_0+v_2$ & $v_{150}$ & $v_{151}$ & $v_{152}$\\
4 & 3 & $v_{102}$ & $v_0$ & $v_{141}$ & $v_{142}$ & $v_{143}$\\
5 & 4 & $v_{106}$ & $v_8-1$ & $v_{153}$ & $v_{154}-1$ & $v_{155}$\\
6 & 4 & $v_{109}$ & $v_9$ & $v_{156}$ & $v_{157}$ & $v_{158}$\\
7 & 4 & $v_{112}$ & $v_{10}$ & $v_{159}$ & $v_{160}$ & $v_{161}$\\
8 & 4 & $v_{115}$ & $v_{11}$ & $v_{162}$ & $v_{163}$ & $v_{164}$\\
9 & 4 & $v_{118}$ & $v_{12}$ & $v_{165}$ & $v_{166}$ & $v_{167}$\\
10 & 4 & $v_{122}$ & $v_{13}$ & $v_{171}$ & $v_{172}$ & $v_{173}$\\
11 & 4 & $v_{125}$ & $v_{14}$ & $v_{174}$ & $v_{175}$ & $v_{176}$\\
12 & 4 & $v_{130}$ & $v_{15}$ & $v_{180}$ & $v_{181}$ & $v_{182}$\\
13 & 3 & $v_{119}$ & $v_4$ & $v_{168}$ & $v_{169}$ & $v_{170}$\\
14 & 3 & $v_{126}$ & $v_1$ & $v_{177}$ & $v_{178}$ & $v_{179}+1$\\
15 & 4 & $v_{133}$ & $2v_{16}$ & $v_{183}$ & $v_{184}$ & $v_{185}$\\
\bottomrule
\end{tabular}
\end{table}


There are sixteen requests and fourteen distinct principal pairs.
Requests sharing the same $(D_i,Y_i,Z_i)$ also have the same $k_i$.
Use one copy of each distinct pair in the following sum:
\[
 \Omega=\sum_{\text{distinct pairs }i}
       \bigl((X_i-1)^2+(Y_i-k_i)^2\bigr).
\]
For $0\leq j\leq4$, let $I_j$ be the set of requests in base group $j$,
and define ordinary polynomial abbreviations
\begin{align}
 d_A&=A-1,& T_j&=2An_j-n_j^2-1,\label{eq:T}\\
 H_j&=\sum_{i\in I_j}Y_i(m_i-X_i),&
 S_j&=\sum_{i\in I_j}\bigl((m_i-D_i)^2-Y_i^2\bigr),\label{eq:HS}\\
 E_j&=S_j-2n_jH_j,&\Pi&=\prod_{j=0}^4T_j.\label{eq:E}
\end{align}
The joint residual is the following polynomial, without a division:
\begin{equation}
 \boxed{\mathcal J=
 \Omega\Pi+d_A\sum_{j=0}^4 n_jE_j\prod_{\ell\ne j}T_\ell
       -z\,d_A\Pi.}
 \label{eq:joint}
\end{equation}
All sums and products have the fixed finite index sets just printed.
They introduce no quantified variables. The intermediate polynomial is
\begin{equation}
 P_{209}=\mathcal J^2+
       \sum_{j\in\{0,\ldots,157\}\setminus D_R}R_j^2.
 \label{eq:P209}
\end{equation}
It has 148 residual squares and degree at most 26.

\section{Why orientation and powers can share one certificate}
\label{sec:jointproof}

This section supplies the unbounded argument, including its order of
dependence. It is essential to classify the raw Pell solutions before
identifying their indices with the requested exponents.

\subsection{Information available before the new residual is used}
The retained signed frame and squared-anchor comparison equations give
\[
 W\geq258,\qquad |A|>W^W,\qquad
 k_i^2,n_j^2,m_i^2,C^2<W.
\]
Set $b=\lfloor\sqrt{W-1}\rfloor$, so $b\geq16$. The comparison
congruences identify the anchor's positive classified index as $|C|$.
The squared-ordinate bound then bounds the absolute classified indices
of all principal pairs by $|C|\leq b$. These conclusions use neither
the removed orientation row nor the removed power rows.

Let $\mathsf X_n(A),\mathsf Y_n(A)$ denote the signed Pell sequences.
Every raw principal pair can therefore be written
\begin{equation}
 (X_i,Y_i)=\epsilon_i\bigl(\mathsf X_{\nu_i}(A),
                              \mathsf Y_{\nu_i}(A)\bigr),
 \quad \epsilon_i\in\{-1,1\},\quad \nu_i\in\Z,
 \quad |\nu_i|\leq b.
 \label{eq:raw}
\end{equation}
At index zero this includes both $(1,0)$ and $(-1,0)$.
For negative $A$, the signed recurrence handles the parity signs; no
positivity of every abscissa is imposed.

\subsection{The small integer with a common sign}
Modulo $d_A=A-1$, the canonical recurrences give
$\mathsf X_\nu(A)\equiv1$ and
$\mathsf Y_\nu(A)\equiv\nu$. Consequently
\begin{equation}
 \Omega\equiv o:=\sum_i\bigl((\epsilon_i-1)^2+
                         (\epsilon_i\nu_i-k_i)^2\bigr)
 \pmod {d_A},\qquad o\geq0.
 \label{eq:o}
\end{equation}
At a nonzero base, $\gcd(n_j,T_j)=1$ because $T_j\equiv-1\pmod{n_j}$.
The Pell recurrence reduces modulo $T_j$ to
\[
 X_i+(n_j-A)Y_i\equiv\epsilon_i n_j^{\nu_i}\pmod{T_j}.
\]
Negative exponents in this proof-only congruence are defined by the
modular inverse of $n_j$. Define
$\delta_i=m_i-X_i+(A-n_j)Y_i=m_i-D_i-n_jY_i$.
A direct polynomial calculation, using $Z_i=AY_i$, gives
\begin{equation}
 \sum_{i\in I_j}\delta_i^2
     =E_j-T_j\sum_{i\in I_j}Y_i^2.
 \label{eq:squareidentity}
\end{equation}

All denominators in the argument are nonzero. The number $d_A$ has the
sign of $A$ and $|d_A|\geq W^W$. If $n_j\ne0$, then
$|T_j|>W^W$ and $n_j/T_j$ has the sign of $A$. If $n_j=0$, then
$T_j=-1$, and the weighted summand vanishes. Dividing
\eqref{eq:joint} and applying \eqref{eq:squareidentity} shows that
\[
 \frac{\Omega}{d_A}+\sum_j\sum_{i\in I_j}
                 \frac{n_j\delta_i^2}{T_j}
\quad\hbox{is an integer.}
\]

Put $C_* =\prod_{j:n_j\ne0}n_j^b$, a nonzero integer. For a request in
a nonzero group define
\[
 e_i=C_*m_i-\epsilon_i\frac{C_*}{n_j^b}n_j^{b+\nu_i}.
\]
Here $C_*/n_j^b$ is an integer and $b+\nu_i\geq0$; thus this
definition uses ordinary integer powers in the proof. We have
$C_*\delta_i\equiv e_i\pmod{T_j}$. Reducing each numerator in its
own denominator proves that
\begin{equation}
 V_*=\frac{C_*^2o}{d_A}+
       \sum_{j:n_j\ne0}\sum_{i\in I_j}\frac{n_je_i^2}{T_j}
 \quad\hbox{is an integer.}
 \label{eq:smallinteger}
\end{equation}
No coprimality between different denominators is required. This is
particularly useful when an index modulus and a power modulus have
large common factors.

Every nonzero summand of \eqref{eq:smallinteger} has the same sign.
The conservative bounds of the joint-certificate theorem accommodate
at most fifteen pairs and seventeen requests; our fourteen and sixteen
are within these bounds:
\begin{align*}
 o&\leq120b^2,& |C_*|&\leq b^{5b},\\
 C_*^2o+\sum_{j,i}|n_j|e_i^2
   &\leq120b^{10b+2}+68b^{12b+1}
     \leq188b^{12b+2}<W^{6b+2}\leq W^W.
\end{align*}
For completeness, $|e_i|\leq2b^{6b}$ follows from $|m_i|\leq b$,
$|C_*|\leq b^{5b}$, and $|n_j^{\nu_i}|\leq b^b$ when
$\nu_i\geq0$; negative indices only decrease that absolute-value
bound. There are at most seventeen terms, each carrying a factor
$|n_j|\leq b$. The strict comparison uses $188<W$, $b^2<W$, and
$6b+2\leq b^2+1\leq W$ for $b\geq16$.

It follows that $|V_*|<1$. Since $V_*$ is an integer, it is zero.
The common-sign property then forces $o=0$ and every $e_i=0$.
Equation~\eqref{eq:o} gives $\epsilon_i=1$ and $\nu_i=k_i$ for
each pair, and $e_i=0$ gives the exact power $m_i=n_j^{k_i}$ at every
nonzero base. When $|n_j|\geq2$, an integer value of this power
excludes a negative $k_i$.

\begin{proposition}[Joint authentication]
Under the retained frame, norm, gate, and squared-anchor hypotheses,
$\mathcal J=0$ simultaneously authenticates the canonical signs and
indices of all fourteen pairs and the sixteen powers at every nonzero
base. Conversely, canonical pairs and exact powers admit an integer
$z$ making $\mathcal J=0$.
\end{proposition}
\begin{proof}
Soundness was proved above. For completeness, the canonical base-one
congruences give $d_A\mid\Omega$. Exact powers and the modular Pell
identity give $T_j\mid\delta_i$, hence $T_j\mid E_j$ by
\eqref{eq:squareidentity}. Therefore
\[
 z=\frac{\Omega}{d_A}+\sum_{j:n_j\ne0}\frac{n_jE_j}{T_j}
\]
is an integer satisfying \eqref{eq:joint}. Zero-base contributions are
zero. These divisions construct witnesses in the proof; none occurs
in the defining polynomial.
\end{proof}

\subsection{Direct restoration, including zero bases}
The new source need not repeat the parent's longer arithmetic decoding
argument. It suffices to restore the eleven deleted source coordinates.
For a nonzero base, joint authentication gives the canonical congruence
and the exact requested power, so $T_j\mid\delta_i$ and $T_j\mid E_j$.
For a zero base, $T_j=-1$, so both divisibilities hold automatically.
Thus the temporary zero-base branches create no restoration gap.

Restore the three removed singleton moduli as $T_j$ and their
quotients as $-\delta_i/T_j$. For either old pooled base group restore
\[
 L_j=E_j/T_j,\qquad K_j=2H_j+(2A-n_j)L_j.
\]
The preceding argument proves these are integers and makes both old
pooled rows zero. Restore the orientation quotient as
$\Omega/d_A$. Source row 147 differs from
$\Omega-d_Av_{204}$ by minus the sum of the fourteen retained
norm rows; this exact identity is included in the independent source
audit. Thus every $P_{209}$ zero extends to a source219 zero. Conversely,
the old source rows themselves establish the divisibilities needed to
construct $z$ in the displayed completeness formula. Every source219
zero therefore projects to a $P_{209}$ zero. No claim of an
integer restoration away from the zero sets is needed here.

\section{Eliminating 124 coordinates over the integers}
\label{sec:elimination}

The preceding reduction is semantic: its quotient witnesses are
restored only after proving divisibility. The next reduction is purely
polynomial and has an integer reconstruction for every free input.

\begin{lemma}[Acyclic unit elimination]
\label{lem:eliminate}
Let an integer polynomial system contain selected rows
$\epsilon_v v+F_v=0$, where $\epsilon_v\in\{-1,1\}$ and $F_v$
does not contain $v$. Suppose the dependence graph among selected
variables, with $v$ depending on those selected variables occurring in
$F_v$, is acyclic. Assign arbitrary integers to the unselected variables
and recursively set $v=-\epsilon_v F_v$. This defines an integer
polynomial map $\Phi$ from the free coordinates to all coordinates.
The remaining equations after substitution have an integer zero if
and only if the original system does. Projection and $\Phi$ are
inverse on the two zero sets.
\end{lemma}
\begin{proof}
Topologically order the selected variables with dependencies first.
At each step the displayed substitution is an integer polynomial in
already defined coordinates. It makes its selected row zero
identically. Any zero of the original system must have exactly these
coordinate values, by induction along the same order. The remaining
rows then vanish on either side precisely when their substituted
versions vanish. This proves both directions and the inverse claim.
\end{proof}

The unit condition is crucial. Solving $2v=u$ over the rationals and
discarding the equation would lose the parity condition on $u$.
Solving $av=F$ without a divisibility certificate can introduce
spurious integer inputs. Our 124 selected coefficients are all
literal $1$ or $-1$, and no selected variable occurs elsewhere in
its defining row.

Appendix~\ref{app:free} lists the 85 free source indices and
Appendix~\ref{app:pivots} gives every selected row, variable, and
unit coefficient in dependency order. Let $\Phi$ be the resulting
map. The 23 unselected old rows are
\begin{equation}
 \begin{split}
 \mathcal K=\{&1,42,83,99,101,109,111,113,115,117,119,121,123,\\
              &125,127,129,131,133,138,140,149,150,151\}.
 \end{split}
 \label{eq:K}
\end{equation}
The principal equation, now fully specified, is
\begin{equation}
 \boxed{P_{85}(\boldsymbol y)=
       \mathcal J(\Phi(\boldsymbol y))^2+
       \sum_{j\in\mathcal K}R_j(\Phi(\boldsymbol y))^2=0.}
 \label{eq:85}
\end{equation}
The fixed finite source row table, request table, free-index list, and
pivot table determine every coefficient of this polynomial. They are
one finite definition, not a program that searches for a suitable
polynomial or assumes an unspecified universal parameter. Source
coordinates outside the free list are polynomial abbreviations, not
additional unknowns.

Since selected rows pull back identically to zero, the exact polynomial
identity $P_{85}=P_{209}\circ\Phi$ also holds. The reverse map on zero
sets is coordinate projection. This makes the existential equivalence
particularly straightforward to formalize in ZFC.

\section{Exact complexity certificates}
\label{sec:metrics}

\subsection{Degree is proved, including noncancellation}
An arithmetic DAG propagates a valid upper degree bound by taking a
maximum at additions and adding degrees at multiplications. Such a
bound alone is not an exact-degree proof: leading terms can cancel.
For the principal equation the retained residual degree bounds are
\[
 (6,4096,8482,50,6,6,6,6,6,6,6,6,6,6,6,6,6,6,
   14,12,2,2,8,39).
\]
Their order is the increasing list $\mathcal K$, followed by
$\mathcal J$. The maximal residual is source row 83, of degree at
most 8,482 after substitution.

For a completely exact lower bound, put $p=1\,000\,000\,007$ and let
$s_0<\cdots<s_{84}$ be the free source indices. Evaluate the propagated
top homogeneous component on the integer direction
$a_i=s_i+3$. Its residue at the row-83 residual is
\[
 [\lambda^{8482}]R_{83}(\Phi(\lambda\boldsymbol a))
       \equiv 140\,621\,492\not\equiv0\pmod p.
\]
The leading term has a simple source-level explanation. In row 83,
$-v_{79}v_{95}$ is the unique highest-degree term after reconstruction.
The reconstructed coordinates have exact degrees 8,198 and 284;
their top components on the same ray are respectively $46\,952\,950$
and $64\,682\,962$ modulo $p$. Every other term of row 83 has degree
at most 12. Their product gives the displayed nonzero residue.
This modular calculation is an exact finite witness that the integer
coefficient is nonzero; it is not a probabilistic polynomial-identity
test. Therefore that residual has exact degree 8,482.
Over the real numbers a sum of polynomial squares cannot cancel its
maximal homogeneous components. Equivalently, on this integer ray,
the leading real coefficient of the sum is a sum of squares with a
nonzero summand. Hence $\degp P_{85}=2\cdot8482=16\,964$.
The last step is over $\mathbb R$ or $\mathbb Q$, not a claim that
sums of squares cannot cancel modulo $p$.

\subsection{Every counted variable actually occurs}
A syntactically live input need not affect the resulting polynomial;
the expression $x-x$ demonstrates the distinction. The archive gives
one exact modular point and all 85 partial derivative values computed
by reverse differentiation of the literal DAG. Every value is nonzero.
Thus each integer polynomial $\partial P_{85}/\partial y_i$ is
nonzero, and all 85 coordinates genuinely occur. The point was found
by a deterministic search, but verification of the supplied witness
uses only exact integer arithmetic modulo the displayed modulus.

\subsection{Coefficient and evaluator accounting}
The checker propagates upper bounds on the coefficient $\ell^1$ norm.
For integer polynomials $f,g$,
\[
 \|f\pm g\|_1\leq\|f\|_1+\|g\|_1,
 \qquad\|fg\|_1\leq\|f\|_1\|g\|_1.
\]
Bounding each norm by a power of two and adding the residual squares
gives $\max_\alpha|[\boldsymbol y^\alpha]P_{85}|\leq2^{59805}$.
No huge expanded polynomial needs to be materialized to prove this
bound. The largest circuit literal $2^{40}+4$ and the expanded-height
bound are deliberately reported separately.

The JSON specification retains some nodes solely for reconstruction
and certificate comparison. The operation count includes only binary
operations reachable from the final polynomial output. The principal
file contains 3,559 stored nodes, while its live evaluator uses 1,788
binary operations. Constants and input nodes are not counted as
binary operations. These conventions match the stated comparison
metric and do not assert optimal arithmetic-circuit size.

\section{Logical calibration and internal ZFC meaning}
\label{sec:logic}

The inherited chain is
\[
 \begin{gathered}
 \text{square-shell C222 refutation}
 \ \longleftrightarrow\ \text{original C222 derivation of code 1},\\
 \text{original C222 derivation of code 1}
 \ \longleftrightarrow\ \text{a finite ZF refutation},\\
 \text{a finite ZF refutation}
 \ \longleftrightarrow\ \text{a finite ZFC refutation}.
 \end{gathered}
\]
The named Lean source interfaces are
\path{KEndpoint.c222_iff_zf_refutation} and
\path{ConstructibleChoice.zf_refutation_iff_zfc_refutation}.
The effective first-order axiom predicate is the repository's
\texttt{ZFCax\_s}, with finite proof relation \texttt{BProv} and
falsity \texttt{Form.fBot}. The square-shell syntax change is an
ordinary effective recoding; it is not a literal identification of
two differently coded formula constructors.

The source calculus-to-ZF bridge is a substantial inherited theorem.
Its model arguments use the internally well-founded part of a
quotient construction, and the ZF-to-ZFC consistency transport uses
the constructible model. None of this is replaced by assuming that
a consistent theory has an externally well-founded or transitive
model. The present packet identifies and retains the relevant source
proofs; it does not freshly compile the full transitive import closure.
The source audit and provenance ledger should be read as such, not
as a new kernel-checked theorem about the entire development.

For the new steps, the formalizability claim is direct. Pell
recurrences, their classification and congruences, finite CRT,
polynomial substitution, prime factorization, and induction over
finite syntactic objects are ordinary constructions provable in ZFC.
The inequalities in Section~\ref{sec:jointproof} are uniform in the
internally coded integers. They do not quantify only over externally
standard proof codes. In a nonstandard model of ZFC, the theorem
concerns that model's integers and finite derivations, interpreted by
the same internally proved lemmas.

Combining the inherited source theorem, the parent-restoration
equivalence, and Lemma~\ref{lem:eliminate} proves
Theorem~\ref{thm:main}. Negating the existential refutation statement
gives exactly the usual arithmetized $\Con(\ZFC)$. It gives no
independent proof that $P_{85}$ has no integer zero.

\section{Verification, reproduction, and integration}
\label{sec:verification}

The archive is self-contained for reproducing the new finite
polynomial checks. Python 3.10 or later and its standard library
suffice; a TeX installation is needed only to rebuild this PDF.
From the extracted package root, run
\begin{verbatim}
python3 -B verify.py
python3 -O -B verify.py
\end{verbatim}
The verifier checks the source hashes, exact retained/deleted row
interface, the joint pool identities, unit pivots, acyclicity,
complete substitution of source DAG nodes, the sum of squares,
exact degree witnesses, all-variable dependence witnesses, and
reported finite metrics. Deliberate corruptions are required to fail.
It imports neither the projection producer nor its search routine.
Optimization-mode replay checks that substantive gates are not
implemented merely as removable Python assertions.

The source and new proofs establish unbounded facts. The finite
verifier authenticates their literal algebraic instances. Passing
checks does not, by itself, establish the Pell classification,
logical calibration, or absence of solutions. No search for a
Diophantine zero is performed or needed.

The principal JSON contains a topologically ordered list of constants,
variables, and binary operations; its final output is the actual
integer polynomial. A supplied evaluator can inspect its input order,
evaluate it exactly, or evaluate it modulo a chosen positive integer.
A modular zero would not constitute an integer solution. Evaluating
the high-degree forms exactly can produce large values even at
small inputs.

For repository intake, add the ZIP as an independent manuscript
package in \texttt{docs/incoming}, then place it according to the
repository's intake procedure. Preserve the source ledger, inherited
notices, proof-boundary statement, and domain distinction. The package
does not replace the maintained fifteen-positive-variable display,
claim an improvement in every simplicity measure, or silently
upgrade an ordinary proof to a fresh Lean/Rocq build.

\section{An unsuccessful route worth recording}
\label{sec:obstruction}

We also examined the latest ten-positive-variable degree-208
construction of package~178 \cite{p178}. Its center approximates a square root
of a product of Pell expressions using polynomial quotients at
infinity. A natural extension is to retain more terms in
\[
 \sqrt{A^2-1}=A-\frac1{2A}-\frac1{8A^3}
                    -\frac1{16A^5}-\cdots.
\]
The degree relevant to the final equation is a weighted total degree,
not just the degree in $A$. Taking the full polynomial part in $A$
does not improve the parent's degree: a term involving high powers
of $Y$ survives with the same weighted degree.

In the actual variables $A=XY+Y+1$, with weights
$(X,Y,k,U,b,c)=(2,2,1,7,4,3)$, the first nonpolynomial term of
the relevant scale-four root expansion is
$3X^{-1}Y^7k^4U^3b$, of weight 41. The next terms are
$-3X^{-2}Y^7k^4U^3b+12X^{-1}Y^6k^4U^3b$, of weight 39.
Multiplication by $X$ permits the first correction but leaves a
defect of weight 108, compared with the parent's 106; its final
degree is worse. Multiplication by $X+1$ can cancel one further
term but leaves the same defect weight 106. The independently
checked notes in \texttt{research/} retain the exact restricted
calculation and its assumptions.

A stronger restricted statement follows from these coefficients.
Let
\[
 \begin{gathered}
 A=XY+Y+1,\quad E=1+XY^2k^2+X^2Y^4k^2,\quad\Delta=A^2-1,\\
 G=E+\Delta U^2E^2,\qquad V=(1+2\Delta U^2E)b-c.
 \end{gathered}
\]
\begin{proposition}[Affine-scale obstruction]
In $\mathbb Q[X,Y,k,U,b,c]$ with the weights above, for every nonzero
$r=aX+b_0Y+d\in\mathbb Q[X,Y]$ and every polynomial $C$,
\[
 \deg_w(C^2-16r^2GV^2)\geq106.
\]
\end{proposition}
\begin{proof}
Introduce a weighted Laurent parameter and choose the formal root
$S=4\sqrt G\,V$ with leading monomial
$8X^7Y^{11}k^4U^3b$ of weight 65. Above weight 41 its components
are polynomials. Write $K=k^4U^3b$. Its nonpolynomial components
at weights 41, 39, and 37 are exactly
\[
 \frac{3KY^7}{X},\qquad
 -\frac{3KY^7}{X^2}+\frac{12KY^6}{X},\qquad
 \frac{3KY^7}{X^3}-\frac{15KY^6}{X^2}+\frac{15KY^5}{X}.
\]
They are obtained by recursively squaring and comparing coefficients;
the supplied standard-library checker verifies these finite identities
with rational coefficients. For nonzero constant $r$, a defect below
106 requires agreement of $C$ with one sign of $rS$ through weight
41, which its first nonpolynomial term prevents. For nonconstant $r$,
the leading weight is 67, so a defect below 106 requires agreement
through weight 39. At weight 43 the nonpolynomial part is
$3b_0KY^8/X$, forcing $b_0=0$. At weight 41 it is
$3(d-a)KY^7/X$, forcing $d=a$. The remaining nonzero scale is
$a(X+1)$, whose weight-39 nonpolynomial part is $-3aKY^6/X$.
This is impossible for a polynomial center. The degree conclusion
uses $(C-rS)(C+rS)$: after choosing the leading sign, the second
factor retains weight 65 or 67, so its product with the first
unmatched coefficient cannot cancel in the rational function field.
\end{proof}

This is a theorem about the displayed abstract polynomial ring and
affine scales. It is not a global lower bound for consistency equations,
arbitrary nonlinear scales, or different encodings. In particular, it
prevents recycling this specific unsuccessful improvement attempt
without claiming that degree 208 is universally minimal.

\section{Further research questions}

\begin{enumerate}
\item Can the joint orientation-and-power residual be applied before
more of the source's quadratic coordinates are introduced? The
present composition saves certificate variables and then projects;
a compiler that preserves this sharing from the outset might reduce
both source size and literal coefficient growth.
\item Can a smaller constructor majorant replace $2^{20}$ while
preserving the boundary-valid multiset-separation proof? This targets
the main coefficient tradeoff rather than just the variable count.
An improved bound must cover every local constructor on every
allowed parameter tuple, not only sampled syntax.
\item Which dependency structures prevent the remaining unit
eliminations? A certified search over integer coordinate shears,
followed by acyclic unit pivots, could expose equations that are
solvable integrally but not monic in the current coordinates.
Any nonunit pivot requires an explicit congruence or divisibility
argument.
\item Can the new low-dimensional forms be given sharper expanded
coefficient bounds, or smaller evaluation circuits, without changing
their zero sets? Circuit simplification should preserve an exact
map to the frozen definition and recount all live constant products.
\item Can the new source-interface theorem and reconstruction map
be formalized on top of the existing ZFC endpoint? A modest first
target is the 209-to-85 polynomial pullback and zero-set equivalence;
the larger target is an explicit object-language ZFC derivation of
the full consistency equivalence.
\item Can a different complete proof calculus improve the entire
description, rather than its arithmetic certificate alone? The
comparison should charge syntax, coefficients, and both directions
of the internal consistency calibration.
\end{enumerate}

\appendix
\clearpage
\section{The 85 free coordinates}
\label{app:free}
Write the following ordered source-index list as
$(s_0,\ldots,s_{84})$. The only independent unknowns of the principal
equation are $y_0,\ldots,y_{84}\in\Z$, with $v_{s_i}=y_i$.
In particular $v_{219}$ is the joint quotient. All other coordinates
that occur below are fixed polynomial abbreviations defined by the
next appendix.
\[
\begin{gathered}
3,5,6,8,9,10,11,12,13,14,15,16,17,18,\\
57,99,100,101,102,103,106,107,109,110,112,113,115,116,\\
118,119,120,122,123,127,128,131,132,133,134,138,141,142,\\
150,151,153,154,156,157,159,160,162,163,165,166,168,169,\\
171,172,174,175,177,178,180,181,183,184,186,187,190,191,\\
196,197,198,199,200,205,206,207,208,209,210,211,212,218,\\
219.
\end{gathered}
\]


\section{The complete integral reconstruction}
\label{app:pivots}
For each table row $(v,j,\epsilon)$, remove the unique monomial
$\epsilon v_v$ from the literal source polynomial $R_j$ in the next
appendix. Call the resulting polynomial $F_v$. Define
$v_v=-\epsilon F_v$, using previously defined coordinates. The table
is in dependency order, not merely in increasing variable order.
Thus these 124 instructions, together with the free list, are a
complete division-free polynomial reconstruction. They are also a
shorter finite specification than expanding every high-degree
residual into monomials.
\begin{longtable}{rrrr@{\qquad}rrrr}
\caption{Reconstruction instructions, read left-to-right and then down.}\label{tab:pivots}\\
\toprule
Step & $v$ & Row $j$ & $\epsilon$ & Step & $v$ & Row $j$ & $\epsilon$\\
\midrule\endfirsthead
\toprule
Step & $v$ & Row $j$ & $\epsilon$ & Step & $v$ & Row $j$ & $\epsilon$\\
\midrule\endhead
\bottomrule\endfoot
1 & 0 & 136 & 1 & 2 & 2 & 3 & 1\\
3 & 4 & 157 & 1 & 4 & 7 & 0 & 1\\
5 & 1 & 2 & 1 & 6 & 19 & 4 & 1\\
7 & 20 & 5 & 1 & 8 & 21 & 6 & 1\\
9 & 22 & 7 & 1 & 10 & 23 & 8 & 1\\
11 & 24 & 9 & 1 & 12 & 25 & 10 & 1\\
13 & 26 & 11 & 1 & 14 & 27 & 12 & 1\\
15 & 28 & 13 & 1 & 16 & 29 & 14 & 1\\
17 & 30 & 15 & 1 & 18 & 31 & 16 & 1\\
19 & 32 & 17 & 1 & 20 & 33 & 18 & 1\\
21 & 34 & 19 & 1 & 22 & 35 & 20 & 1\\
23 & 36 & 21 & 1 & 24 & 37 & 22 & 1\\
25 & 38 & 23 & 1 & 26 & 39 & 24 & 1\\
27 & 40 & 25 & 1 & 28 & 41 & 26 & 1\\
29 & 42 & 27 & 1 & 30 & 43 & 28 & 1\\
31 & 44 & 29 & 1 & 32 & 45 & 30 & 1\\
33 & 46 & 31 & 1 & 34 & 47 & 32 & 1\\
35 & 48 & 33 & 1 & 36 & 49 & 34 & 1\\
37 & 50 & 35 & 1 & 38 & 51 & 36 & 1\\
39 & 52 & 37 & 1 & 40 & 53 & 38 & 1\\
41 & 54 & 39 & 1 & 42 & 55 & 40 & 1\\
43 & 56 & 41 & 1 & 44 & 58 & 43 & 1\\
45 & 59 & 44 & 1 & 46 & 60 & 45 & 1\\
47 & 61 & 46 & 1 & 48 & 62 & 47 & 1\\
49 & 63 & 48 & 1 & 50 & 64 & 49 & 1\\
51 & 65 & 50 & 1 & 52 & 66 & 51 & 1\\
53 & 67 & 52 & 1 & 54 & 68 & 53 & 1\\
55 & 69 & 54 & -1 & 56 & 70 & 55 & -1\\
57 & 71 & 56 & -1 & 58 & 72 & 57 & -1\\
59 & 73 & 58 & -1 & 60 & 74 & 59 & -1\\
61 & 75 & 60 & -1 & 62 & 76 & 61 & -1\\
63 & 77 & 62 & 1 & 64 & 78 & 63 & 1\\
65 & 80 & 65 & 1 & 66 & 79 & 64 & 1\\
67 & 81 & 67 & 1 & 68 & 213 & 152 & 1\\
69 & 82 & 68 & 1 & 70 & 83 & 66 & 1\\
71 & 84 & 69 & 1 & 72 & 85 & 70 & 1\\
73 & 86 & 71 & 1 & 74 & 87 & 72 & 1\\
75 & 88 & 73 & 1 & 76 & 89 & 74 & 1\\
77 & 90 & 75 & 1 & 78 & 91 & 76 & 1\\
79 & 92 & 77 & 1 & 80 & 93 & 78 & 1\\
81 & 94 & 79 & 1 & 82 & 95 & 80 & 1\\
83 & 96 & 81 & 1 & 84 & 97 & 82 & 1\\
85 & 98 & 84 & 1 & 86 & 126 & 148 & 1\\
87 & 105 & 85 & 1 & 88 & 104 & 86 & -1\\
89 & 129 & 93 & 1 & 90 & 108 & 87 & -1\\
91 & 111 & 88 & -1 & 92 & 114 & 89 & -1\\
93 & 117 & 90 & -1 & 94 & 121 & 91 & -1\\
95 & 124 & 92 & -1 & 96 & 125 & 94 & 1\\
97 & 130 & 95 & 1 & 98 & 135 & 96 & -1\\
99 & 189 & 135 & 1 & 100 & 136 & 134 & 1\\
101 & 203 & 141 & 1 & 102 & 192 & 142 & 1\\
103 & 195 & 143 & 1 & 104 & 137 & 144 & -1\\
105 & 139 & 97 & 1 & 106 & 140 & 98 & 1\\
107 & 143 & 100 & 1 & 108 & 152 & 108 & 1\\
109 & 155 & 110 & 1 & 110 & 158 & 112 & 1\\
111 & 161 & 114 & 1 & 112 & 164 & 116 & 1\\
113 & 167 & 118 & 1 & 114 & 170 & 120 & 1\\
115 & 173 & 122 & 1 & 116 & 176 & 124 & 1\\
117 & 179 & 126 & 1 & 118 & 182 & 128 & 1\\
119 & 185 & 130 & 1 & 120 & 188 & 132 & 1\\
121 & 193 & 145 & 1 & 122 & 194 & 146 & 1\\
123 & 201 & 137 & 1 & 124 & 202 & 139 & 1\\
\end{longtable}


\clearpage
\section{Complete retained source row table}
\label{app:source}
The following 147 quadratic polynomials are the retained rows from
the literal source219 system. Source labels use zero-based indices.
Rows listed in the preceding pivot table define coordinates; the 23
rows in \eqref{eq:K} contribute residual squares in \eqref{eq:85}.
The additional joint residual is completely defined by
\eqref{eq:T}--\eqref{eq:joint} and Table~\ref{tab:requests}.
No other equations or unspecified constants enter the definition.

\begingroup
\small
\allowdisplaybreaks[4]
\setlength{\jot}{3.8pt}
\begin{align*}
R_{0}&=v_{0}-v_{0}v_{4}-v_{4}+v_{7}\\[2pt]
R_{1}&=-1+v_{1}-v_{2}-v_{4}v_{5}\\[2pt]
R_{2}&=-1+v_{1}-v_{6}v_{7}\\[2pt]
R_{3}&=1+v_{2}-v_{3}v_{18}+2v_{16}\\[2pt]
R_{4}&=-1+2v_{8}-v_{8}^{2}-2v_{8}v_{9}+2v_{9}-v_{9}^{2}+v_{19}\\[2pt]
R_{5}&=-5-30v_{8}-36v_{8}^{2}-96v_{8}v_{9}-24v_{8}v_{19}-40v_{9}-64v_{9}^{2}-32v_{9}v_{19}\\*
&\quad-10v_{19}-4v_{19}^{2}+v_{20}\\[2pt]
R_{6}&=-2+2v_{8}-v_{8}^{2}-4v_{8}v_{20}+4v_{20}-4v_{20}^{2}+v_{21}\\[2pt]
R_{7}&=1+6v_{8}-36v_{8}^{2}-192v_{8}v_{20}-24v_{8}v_{21}+16v_{20}-256v_{20}^{2}-64v_{20}v_{21}\\*
&\quad+2v_{21}-4v_{21}^{2}+v_{22}\\[2pt]
R_{8}&=-2-v_{9}^{2}-2v_{9}v_{10}-v_{10}^{2}+v_{23}\\[2pt]
R_{9}&=-9+7v_{8}-v_{8}^{2}-4v_{8}v_{9}-4v_{8}v_{23}+12v_{9}-4v_{9}^{2}-8v_{9}v_{23}+12v_{23}\\*
&\quad-4v_{23}^{2}+v_{24}\\[2pt]
R_{10}&=-1+2v_{8}-v_{8}^{2}-2v_{8}v_{10}+2v_{10}-v_{10}^{2}+v_{25}\\[2pt]
R_{11}&=-6v_{8}-144v_{8}^{2}-192v_{8}v_{9}-192v_{8}v_{10}-48v_{8}v_{19}-48v_{8}v_{25}-8v_{9}\\*
&\quad-64v_{9}^{2}-128v_{9}v_{10}-32v_{9}v_{19}-32v_{9}v_{25}-64v_{10}^{2}-32v_{10}v_{19}\\*
&\quad-32v_{10}v_{25}-2v_{19}-4v_{19}^{2}-8v_{19}v_{25}-4v_{25}^{2}+v_{26}\\[2pt]
R_{12}&=-4v_{24}^{2}-8v_{24}v_{26}-4v_{26}^{2}+v_{27}\\[2pt]
R_{13}&=-6+16v_{10}-16v_{10}^{2}-16v_{10}v_{23}-16v_{10}v_{25}+8v_{23}-4v_{23}^{2}-8v_{23}v_{25}\\*
&\quad+8v_{25}-4v_{25}^{2}+v_{28}\\[2pt]
R_{14}&=1+6v_{8}-36v_{8}^{2}-96v_{8}v_{9}-48v_{8}v_{10}-24v_{8}v_{19}-48v_{8}v_{25}\\*
&\quad-24v_{8}v_{28}+8v_{9}-64v_{9}^{2}-64v_{9}v_{10}-32v_{9}v_{19}-64v_{9}v_{25}\\*
&\quad-32v_{9}v_{28}+8v_{10}-16v_{10}^{2}-16v_{10}v_{19}-32v_{10}v_{25}-16v_{10}v_{28}+2v_{19}\\*
&\quad-4v_{19}^{2}-16v_{19}v_{25}-8v_{19}v_{28}+8v_{25}-16v_{25}^{2}-16v_{25}v_{28}+4v_{28}\\*
&\quad-4v_{28}^{2}+v_{29}\\[2pt]
R_{15}&=-10-5v_{8}-v_{8}^{2}-12v_{8}v_{9}-16v_{8}v_{10}-4v_{8}v_{23}-36v_{9}-36v_{9}^{2}\\*
&\quad-96v_{9}v_{10}-24v_{9}v_{23}-48v_{10}-64v_{10}^{2}-32v_{10}v_{23}-12v_{23}-4v_{23}^{2}\\*
&\quad+v_{30}\\[2pt]
R_{16}&=-2-24v_{8}-36v_{8}^{2}-12v_{8}v_{9}-96v_{8}v_{10}-24v_{8}v_{25}-3v_{9}-v_{9}^{2}\\*
&\quad-16v_{9}v_{10}-4v_{9}v_{25}-32v_{10}-64v_{10}^{2}-32v_{10}v_{25}-8v_{25}-4v_{25}^{2}+v_{31}\\[2pt]
R_{17}&=2-4v_{30}^{2}-8v_{30}v_{31}-4v_{31}^{2}+v_{32}\\[2pt]
R_{18}&=-12+16v_{8}-4v_{8}^{2}-8v_{8}v_{10}-8v_{8}v_{19}-8v_{8}v_{23}+16v_{10}-4v_{10}^{2}\\*
&\quad-8v_{10}v_{19}-8v_{10}v_{23}+16v_{19}-4v_{19}^{2}-8v_{19}v_{23}+16v_{23}-4v_{23}^{2}+v_{33}\\[2pt]
R_{19}&=-3-24v_{8}-64v_{8}^{2}-160v_{8}v_{10}-64v_{8}v_{23}-32v_{8}v_{25}-32v_{8}v_{33}\\*
&\quad-26v_{10}-100v_{10}^{2}-80v_{10}v_{23}-40v_{10}v_{25}-40v_{10}v_{33}-8v_{23}-16v_{23}^{2}\\*
&\quad-16v_{23}v_{25}-16v_{23}v_{33}-6v_{25}-4v_{25}^{2}-8v_{25}v_{33}-4v_{33}-4v_{33}^{2}+v_{34}\\[2pt]
R_{20}&=-1-18v_{8}-36v_{8}^{2}-96v_{8}v_{9}-24v_{8}v_{19}-24v_{9}-64v_{9}^{2}-32v_{9}v_{19}\\*
&\quad-6v_{19}-4v_{19}^{2}+v_{35}\\[2pt]
R_{21}&=-2+2v_{8}-v_{8}^{2}-4v_{8}v_{35}+4v_{35}-4v_{35}^{2}+v_{36}\\[2pt]
R_{22}&=1+6v_{8}-36v_{8}^{2}-192v_{8}v_{35}-24v_{8}v_{36}+16v_{35}-256v_{35}^{2}-64v_{35}v_{36}\\*
&\quad+2v_{36}-4v_{36}^{2}+v_{37}\\[2pt]
R_{23}&=-v_{9}^{2}-4v_{9}v_{37}-4v_{37}^{2}+v_{38}\\[2pt]
R_{24}&=-6-30v_{8}-64v_{8}^{2}-128v_{8}v_{9}-64v_{8}v_{19}-32v_{9}-64v_{9}^{2}-64v_{9}v_{19}\\*
&\quad-14v_{19}-16v_{19}^{2}+v_{39}\\[2pt]
R_{25}&=-4+8v_{9}-64v_{9}^{2}-96v_{9}v_{10}-32v_{9}v_{23}-32v_{9}v_{39}+6v_{10}-36v_{10}^{2}\\*
&\quad-24v_{10}v_{23}-24v_{10}v_{39}+2v_{23}-4v_{23}^{2}-8v_{23}v_{39}+4v_{39}-4v_{39}^{2}+v_{40}\\[2pt]
R_{26}&=8v_{8}-64v_{8}^{2}-96v_{8}v_{10}-32v_{8}v_{25}-32v_{8}v_{40}+6v_{10}-36v_{10}^{2}\\*
&\quad-24v_{10}v_{25}-24v_{10}v_{40}+2v_{25}-4v_{25}^{2}-8v_{25}v_{40}+4v_{40}-4v_{40}^{2}+v_{41}\\[2pt]
R_{27}&=-v_{10}^{2}-2v_{10}v_{11}-v_{11}^{2}+v_{42}\\[2pt]
R_{28}&=2v_{8}-v_{8}^{2}-4v_{8}v_{10}-4v_{8}v_{42}+4v_{10}-4v_{10}^{2}-8v_{10}v_{42}+4v_{42}\\*
&\quad-4v_{42}^{2}+v_{43}\\[2pt]
R_{29}&=-35+24v_{8}-4v_{8}^{2}-8v_{8}v_{9}-8v_{8}v_{23}-8v_{8}v_{25}+24v_{9}-4v_{9}^{2}\\*
&\quad-8v_{9}v_{23}-8v_{9}v_{25}+24v_{23}-4v_{23}^{2}-8v_{23}v_{25}+24v_{25}-4v_{25}^{2}+v_{44}\\[2pt]
R_{30}&=-16v_{8}^{2}-32v_{8}v_{10}-32v_{8}v_{25}-32v_{8}v_{42}-16v_{8}v_{43}-16v_{8}v_{44}\\*
&\quad-4v_{10}-16v_{10}^{2}-32v_{10}v_{25}-32v_{10}v_{42}-16v_{10}v_{43}-16v_{10}v_{44}\\*
&\quad-16v_{25}^{2}-32v_{25}v_{42}-16v_{25}v_{43}-16v_{25}v_{44}-4v_{42}-16v_{42}^{2}\\*
&\quad-16v_{42}v_{43}-16v_{42}v_{44}-2v_{43}-4v_{43}^{2}-8v_{43}v_{44}-4v_{44}^{2}+v_{45}\\[2pt]
R_{31}&=v_{10}-25v_{10}^{2}-40v_{10}v_{25}-20v_{10}v_{28}-16v_{25}^{2}-16v_{25}v_{28}-4v_{28}^{2}\\*
&\quad+v_{46}\\[2pt]
R_{32}&=-1+8v_{9}-4v_{9}^{2}-8v_{9}v_{23}+8v_{23}-4v_{23}^{2}+v_{47}\\[2pt]
R_{33}&=-1-16v_{9}^{2}-8v_{9}v_{10}-32v_{9}v_{23}-16v_{9}v_{47}-v_{10}^{2}-8v_{10}v_{23}\\*
&\quad-4v_{10}v_{47}-16v_{23}^{2}-16v_{23}v_{47}-4v_{47}^{2}+v_{48}\\[2pt]
R_{34}&=1+v_{9}-v_{9}^{2}-4v_{9}v_{10}-4v_{9}v_{48}-4v_{10}^{2}-8v_{10}v_{48}-4v_{48}^{2}+v_{49}\\[2pt]
R_{35}&=-1-4v_{49}-4v_{49}^{2}+v_{50}\\[2pt]
R_{36}&=-1-v_{10}^{2}-8v_{10}v_{49}-4v_{10}v_{50}-16v_{49}^{2}-16v_{49}v_{50}-4v_{50}^{2}+v_{51}\\[2pt]
R_{37}&=2+v_{9}-v_{9}^{2}-4v_{9}v_{10}-4v_{9}v_{51}-4v_{10}^{2}-8v_{10}v_{51}-4v_{51}^{2}+v_{52}\\[2pt]
R_{38}&=-4v_{46}^{2}-8v_{46}v_{52}-4v_{52}^{2}+v_{53}\\[2pt]
R_{39}&=-9+7v_{8}-v_{8}^{2}-8v_{8}v_{46}-4v_{8}v_{53}+24v_{46}-16v_{46}^{2}-16v_{46}v_{53}+12v_{53}\\*
&\quad-4v_{53}^{2}+v_{54}\\[2pt]
R_{40}&=-1-4v_{54}-4v_{54}^{2}+v_{55}\\[2pt]
R_{41}&=-v_{10}^{2}-8v_{10}v_{54}-4v_{10}v_{55}-16v_{54}^{2}-16v_{54}v_{55}-4v_{55}^{2}+v_{56}\\[2pt]
R_{42}&=-v_{8}-v_{8}^{2}-4v_{8}v_{10}-4v_{8}v_{56}-4v_{10}-4v_{10}^{2}-8v_{10}v_{56}-4v_{56}\\*
&\quad-4v_{56}^{2}+v_{57}\\[2pt]
R_{43}&=-10-12v_{57}-4v_{57}^{2}+v_{58}\\[2pt]
R_{44}&=1+4v_{9}-4v_{9}^{2}-8v_{9}v_{23}-16v_{9}v_{57}-8v_{9}v_{58}+4v_{23}-4v_{23}^{2}\\*
&\quad-16v_{23}v_{57}-8v_{23}v_{58}+8v_{57}-16v_{57}^{2}-16v_{57}v_{58}+4v_{58}-4v_{58}^{2}+v_{59}\\[2pt]
R_{45}&=v_{9}-25v_{9}^{2}-40v_{9}v_{23}-20v_{9}v_{59}-16v_{23}^{2}-16v_{23}v_{59}-4v_{59}^{2}+v_{60}\\[2pt]
R_{46}&=-10-12v_{52}-4v_{52}^{2}+v_{61}\\[2pt]
R_{47}&=1-16v_{52}^{2}-16v_{52}v_{60}-16v_{52}v_{61}-4v_{60}^{2}-8v_{60}v_{61}-4v_{61}^{2}+v_{62}\\[2pt]
R_{48}&=v_{10}-v_{10}^{2}-8v_{10}v_{60}-4v_{10}v_{62}-16v_{60}^{2}-16v_{60}v_{62}-4v_{62}^{2}+v_{63}\\[2pt]
R_{49}&=-4v_{8}-4v_{8}^{2}-4v_{8}v_{10}-8v_{8}v_{19}-v_{10}-v_{10}^{2}-4v_{10}v_{19}-4v_{19}\\*
&\quad-4v_{19}^{2}+v_{64}\\[2pt]
R_{50}&=1+4v_{8}-4v_{8}^{2}-8v_{8}v_{19}-8v_{8}v_{64}+4v_{19}-4v_{19}^{2}-8v_{19}v_{64}+4v_{64}\\*
&\quad-4v_{64}^{2}+v_{65}\\[2pt]
R_{51}&=-2-4v_{8}-16v_{8}^{2}-32v_{8}v_{19}-16v_{8}v_{63}-16v_{8}v_{65}-4v_{19}-16v_{19}^{2}\\*
&\quad-16v_{19}v_{63}-16v_{19}v_{65}-4v_{63}^{2}-8v_{63}v_{65}-2v_{65}-4v_{65}^{2}+v_{66}\\[2pt]
R_{52}&=-2-4v_{66}-4v_{66}^{2}+v_{67}\\[2pt]
R_{53}&=-v_{2}v_{13}+v_{68}\\[2pt]
R_{54}&=v_{12}-v_{13}-2v_{22}v_{68}-v_{69}\\[2pt]
R_{55}&=v_{12}-v_{13}-4v_{24}v_{68}-2v_{27}v_{68}-v_{70}\\[2pt]
R_{56}&=v_{12}-v_{13}-2v_{29}v_{68}-6v_{68}-v_{71}\\[2pt]
R_{57}&=v_{12}-v_{13}-4v_{30}v_{68}-2v_{32}v_{68}-v_{72}\\[2pt]
R_{58}&=v_{12}-v_{13}-2v_{34}v_{68}+6v_{68}-v_{73}\\[2pt]
R_{59}&=v_{12}-v_{13}-4v_{37}v_{68}-2v_{38}v_{68}-v_{74}\\[2pt]
R_{60}&=v_{12}-v_{13}-2v_{41}v_{68}-v_{75}\\[2pt]
R_{61}&=v_{12}-v_{13}-2v_{45}v_{68}-4v_{68}-v_{76}\\[2pt]
R_{62}&=-v_{8}v_{9}+v_{8}v_{10}+v_{9}+v_{9}v_{10}-v_{10}-v_{10}^{2}+v_{77}\\[2pt]
R_{63}&=1+2v_{11}+v_{11}^{2}-v_{77}^{2}+v_{78}\\[2pt]
R_{64}&=-v_{12}^{2}+2v_{12}v_{13}+2v_{12}v_{80}-v_{13}^{2}-2v_{13}v_{80}-v_{78}^{2}+v_{79}-v_{80}^{2}\\[2pt]
R_{65}&=-4v_{66}v_{68}-2v_{67}v_{68}+2v_{68}+v_{80}\\[2pt]
R_{66}&=-v_{9}v_{68}-v_{13}+v_{81}v_{82}+v_{83}\\[2pt]
R_{67}&=-1+2v_{2}-2v_{2}v_{8}-2v_{2}v_{19}+v_{81}\\[2pt]
R_{68}&=-v_{8}v_{213}-v_{12}+v_{82}+v_{213}\\[2pt]
R_{69}&=-6v_{8}v_{68}-8v_{9}v_{68}+v_{9}v_{213}+v_{12}-v_{13}-2v_{19}v_{68}+v_{84}\\[2pt]
R_{70}&=-v_{12}v_{69}+v_{13}v_{69}+v_{85}\\[2pt]
R_{71}&=-v_{70}v_{85}+v_{86}\\[2pt]
R_{72}&=-v_{71}v_{86}+v_{87}\\[2pt]
R_{73}&=-v_{72}v_{87}+v_{88}\\[2pt]
R_{74}&=-v_{73}v_{88}+v_{89}\\[2pt]
R_{75}&=-v_{74}v_{89}+v_{90}\\[2pt]
R_{76}&=-v_{75}v_{90}+v_{91}\\[2pt]
R_{77}&=-v_{76}v_{91}+v_{92}\\[2pt]
R_{78}&=-v_{83}v_{92}+v_{93}\\[2pt]
R_{79}&=-v_{84}v_{93}+v_{94}\\[2pt]
R_{80}&=-v_{94}^{2}+v_{95}\\[2pt]
R_{81}&=-v_{4}-2v_{4}v_{16}+2v_{16}+v_{96}\\[2pt]
R_{82}&=-v_{14}v_{96}+v_{97}\\[2pt]
R_{83}&=-2v_{1}^{2}+2v_{1}v_{12}+2v_{1}v_{13}+2v_{1}v_{97}+2v_{1}v_{98}-v_{12}^{2}-2v_{12}v_{97}\\*
&\quad-v_{13}^{2}-2v_{13}v_{98}+v_{17}v_{99}-v_{79}v_{95}-v_{97}^{2}-v_{98}^{2}\\[2pt]
R_{84}&=-v_{4}v_{15}+v_{15}-v_{15}v_{96}+v_{98}\\[2pt]
R_{85}&=-v_{100}v_{103}+v_{103}+v_{105}\\[2pt]
R_{86}&=-v_{3}v_{126}+v_{101}v_{205}-v_{104}-v_{105}v_{126}\\[2pt]
R_{87}&=v_{106}v_{205}-v_{107}v_{129}-v_{108}\\[2pt]
R_{88}&=v_{109}v_{205}-v_{110}v_{129}-v_{111}\\[2pt]
R_{89}&=v_{112}v_{205}-v_{113}v_{129}-v_{114}\\[2pt]
R_{90}&=v_{115}v_{205}-v_{116}v_{129}-v_{117}\\[2pt]
R_{91}&=v_{118}v_{206}-v_{120}v_{129}-v_{121}\\[2pt]
R_{92}&=v_{122}v_{206}-v_{123}v_{129}-v_{124}\\[2pt]
R_{93}&=-v_{17}v_{126}+v_{129}\\[2pt]
R_{94}&=v_{125}-v_{127}v_{129}-v_{128}\\[2pt]
R_{95}&=-v_{129}v_{131}+v_{130}-v_{132}\\[2pt]
R_{96}&=-v_{129}v_{134}+v_{133}v_{205}-v_{135}\\[2pt]
R_{97}&=v_{136}-2v_{136}v_{138}+v_{139}\\[2pt]
R_{98}&=-v_{136}v_{139}-v_{139}+v_{140}\\[2pt]
R_{99}&=2v_{137}+v_{137}^{2}+v_{139}^{2}-v_{140}^{2}\\[2pt]
R_{100}&=-v_{137}v_{142}-v_{142}+v_{143}\\[2pt]
R_{101}&=-1+v_{141}^{2}+2v_{141}v_{143}+v_{142}^{2}\\[2pt]
R_{108}&=-v_{137}v_{151}-v_{151}+v_{152}\\[2pt]
R_{109}&=-1+v_{150}^{2}+2v_{150}v_{152}+v_{151}^{2}\\[2pt]
R_{110}&=1+v_{137}-v_{137}v_{154}-v_{154}+v_{155}\\[2pt]
R_{111}&=v_{153}^{2}+2v_{153}v_{155}-2v_{154}+v_{154}^{2}\\[2pt]
R_{112}&=-v_{137}v_{157}-v_{157}+v_{158}\\[2pt]
R_{113}&=-1+v_{156}^{2}+2v_{156}v_{158}+v_{157}^{2}\\[2pt]
R_{114}&=-v_{137}v_{160}-v_{160}+v_{161}\\[2pt]
R_{115}&=-1+v_{159}^{2}+2v_{159}v_{161}+v_{160}^{2}\\[2pt]
R_{116}&=-v_{137}v_{163}-v_{163}+v_{164}\\[2pt]
R_{117}&=-1+v_{162}^{2}+2v_{162}v_{164}+v_{163}^{2}\\[2pt]
R_{118}&=-v_{137}v_{166}-v_{166}+v_{167}\\[2pt]
R_{119}&=-1+v_{165}^{2}+2v_{165}v_{167}+v_{166}^{2}\\[2pt]
R_{120}&=-v_{137}v_{169}-v_{169}+v_{170}\\[2pt]
R_{121}&=-1+v_{168}^{2}+2v_{168}v_{170}+v_{169}^{2}\\[2pt]
R_{122}&=-v_{137}v_{172}-v_{172}+v_{173}\\[2pt]
R_{123}&=-1+v_{171}^{2}+2v_{171}v_{173}+v_{172}^{2}\\[2pt]
R_{124}&=-v_{137}v_{175}-v_{175}+v_{176}\\[2pt]
R_{125}&=-1+v_{174}^{2}+2v_{174}v_{176}+v_{175}^{2}\\[2pt]
R_{126}&=1-v_{137}v_{178}-v_{178}+v_{179}\\[2pt]
R_{127}&=-1+2v_{177}+v_{177}^{2}+2v_{177}v_{179}+v_{178}^{2}\\[2pt]
R_{128}&=-v_{137}v_{181}-v_{181}+v_{182}\\[2pt]
R_{129}&=-1+v_{180}^{2}+2v_{180}v_{182}+v_{181}^{2}\\[2pt]
R_{130}&=-v_{137}v_{184}-v_{184}+v_{185}\\[2pt]
R_{131}&=-1+v_{183}^{2}+2v_{183}v_{185}+v_{184}^{2}\\[2pt]
R_{132}&=-v_{137}v_{187}-v_{187}+v_{188}\\[2pt]
R_{133}&=-1+v_{186}^{2}+2v_{186}v_{188}+v_{187}^{2}\\[2pt]
R_{134}&=-5-2v_{0}-1099511627780v_{0}^{2}-2v_{0}v_{2}-v_{1}^{2}-2v_{2}^{2}-v_{3}^{2}-v_{4}^{2}+2v_{8}\\*
&\quad-v_{8}^{2}-v_{9}^{2}-v_{10}^{2}-v_{11}^{2}-v_{12}^{2}-v_{13}^{2}-v_{14}^{2}-v_{15}^{2}-4v_{16}^{2}-v_{17}^{2}\\*
&\quad+2v_{100}-2v_{100}^{2}-v_{101}^{2}-v_{102}^{2}-v_{106}^{2}-v_{109}^{2}-v_{112}^{2}-v_{115}^{2}-v_{118}^{2}-v_{119}^{2}\\*
&\quad-v_{122}^{2}-v_{125}^{2}-v_{126}^{2}-v_{130}^{2}-v_{133}^{2}+v_{136}-v_{189}^{2}\\[2pt]
R_{135}&=2v_{3}-4v_{3}v_{8}-4v_{3}v_{9}-4v_{3}v_{10}-4v_{3}v_{11}-4v_{3}v_{12}-4v_{3}v_{13}\\*
&\quad-4v_{3}v_{14}-4v_{3}v_{15}-8v_{3}v_{16}+v_{189}\\[2pt]
R_{136}&=-16+v_{0}-4v_{190}^{2}\\[2pt]
R_{137}&=-v_{137}v_{192}-v_{192}+v_{201}\\[2pt]
R_{138}&=-1+v_{191}^{2}+v_{192}^{2}-v_{201}^{2}\\[2pt]
R_{139}&=-v_{194}-v_{194}v_{195}+v_{202}\\[2pt]
R_{140}&=-1+v_{193}^{2}+v_{194}^{2}-v_{202}^{2}\\[2pt]
R_{141}&=-v_{184}^{2}+v_{203}\\[2pt]
R_{142}&=v_{192}-2v_{197}v_{203}-v_{203}\\[2pt]
R_{143}&=-4v_{192}-8v_{192}v_{196}+v_{195}\\[2pt]
R_{144}&=-v_{137}-v_{191}v_{198}+v_{195}\\[2pt]
R_{145}&=-v_{183}-v_{185}-v_{191}v_{199}+v_{193}\\[2pt]
R_{146}&=-2v_{16}-4v_{184}v_{200}+v_{194}\\[2pt]
R_{148}&=-v_{102}v_{205}+v_{126}\\[2pt]
R_{149}&=-v_{119}v_{206}+v_{126}\\[2pt]
R_{150}&=-2-4v_{16}^{2}-v_{142}^{2}-v_{151}^{2}+2v_{154}-v_{154}^{2}-v_{157}^{2}-v_{160}^{2}-v_{163}^{2}-v_{166}^{2}\\*
&\quad-v_{169}^{2}-v_{172}^{2}-v_{175}^{2}-v_{178}^{2}-v_{181}^{2}+v_{184}^{2}-v_{187}^{2}-v_{207}^{2}-v_{208}^{2}-v_{209}\\*
&\quad-v_{209}^{2}\\[2pt]
R_{151}&=-1-4v_{104}^{2}-4v_{108}^{2}-4v_{111}^{2}-4v_{114}^{2}-4v_{117}^{2}-4v_{121}^{2}-4v_{124}^{2}+v_{126}^{2}\\*
&\quad-4v_{128}^{2}-4v_{132}^{2}-4v_{135}^{2}-v_{210}^{2}-v_{211}^{2}-v_{212}-v_{212}^{2}\\[2pt]
R_{152}&=-v_{2}v_{12}+v_{213}\\[2pt]
R_{157}&=-1-v_{2}-2v_{2}v_{218}+v_{4}\\[2pt]
\end{align*}
\endgroup


\begin{thebibliography}{99}
\raggedright
\bibitem{p173} ProveIt, package 173, \emph{A reused CRT index: an explicit
Diophantine consistency equation}. Source, ordinary proofs, and exact
equations at the comparison snapshot.\newline
\url{https://github.com/VladimirReshetnikov/ProveIt/tree/ecb9a4b4dda617512952525d50f17343700aac52/Computability/DiophantineConsistency/con-zfc-diophantine-forms/173-crt-index-anchor-215}
\bibitem{p176} ProveIt, package 176, \emph{Oriented pooling}, especially
\texttt{proofs/JOINT\_POOLING.md} and
\texttt{proofs/ORIENTED\_POOLING.md}.\newline
\url{https://github.com/VladimirReshetnikov/ProveIt/tree/ecb9a4b4dda617512952525d50f17343700aac52/Computability/DiophantineConsistency/con-zfc-diophantine-forms/176-oriented-pooling}
\bibitem{p177} ProveIt, package 177, \emph{Oriented power pooling and
polynomial root centering}. The stated direct-family comparison.\newline
\url{https://github.com/VladimirReshetnikov/ProveIt/tree/ecb9a4b4dda617512952525d50f17343700aac52/Computability/DiophantineConsistency/con-zfc-diophantine-forms/177-oriented-powers}
\bibitem{p178} ProveIt, package 178, \emph{Polynomial centers for
Diophantine consistency}.\newline
\url{https://github.com/VladimirReshetnikov/ProveIt/tree/ecb9a4b4dda617512952525d50f17343700aac52/Computability/DiophantineConsistency/con-zfc-diophantine-forms/178-polynomial-centers}
\bibitem{p152} ProveIt, package 152, \emph{Lucas corrections}.
Earlier applications of acyclic unit elimination to direct equations.\newline
\url{https://github.com/VladimirReshetnikov/ProveIt/tree/ecb9a4b4dda617512952525d50f17343700aac52/Computability/DiophantineConsistency/con-zfc-diophantine-forms/152-lucas-corrections}
\bibitem{bayer2026} J.~Bayer, M.~David, M.~Ha\ss ler, Y.~Matiyasevich,
and D.~Schleicher, \emph{Diophantine equations over $\Z$: universal
bounds and parallel formalization}, arXiv:2506.20909v3, 22 October 2026.\newline
\url{https://arxiv.org/abs/2506.20909}
\end{thebibliography}
\end{document}
